% II. Background and related work. Every key cited below is in refs.bib and
% was verified (docs/REFERENCES.md). Uses the corrected titles and venues for
% CAPSULE, Savva and Herodotou, Akgun et al. and the ECML PKDD chapter; the
% NVMW'22 abstract is non-archival and is cited as such.

\section{Background and Related Work}\label{sec:related}

\textbf{Read-ahead in the kernel.} Linux keeps a per-file read-ahead window
that grows while the reader stays sequential and collapses on a random
access; a marker page triggers the next asynchronous
batch~\cite{wu2007readahead,wu2010concurrent,linux_mmapi_readahead}, and
applications can add hints~\cite{man_posix_fadvise,man_readahead}. Our
\texttt{linux\_ra} emulates this window; the testbed measures the kernel's
own.

\textbf{Prefetching in storage caches.} Sequential prefetching with adaptive
degree and trigger distance is standard in storage
controllers~\cite{gill2005sarc,gill2007amp}; other work exploits disk layout
and access history~\cite{ding2007diskseen}, mined block
correlations~\cite{li2004cminer} or server-side
context~\cite{soundararajan2008context}. Informed prefetching uses application
hints and bounds how far ahead to fetch~\cite{patterson1995informed}, and
integrated prefetching and caching studies when a prefetch should replace a
cached block~\cite{cao1995integrated}. Predictive file prefetching goes back to
probability graphs and Markov predictors~\cite{griffioen1994reducing,
kroeger2001design,joseph1997markov}.

\textbf{Learned prefetching.} LSTM models of address
deltas~\cite{hochreiter1997lstm,hashemi2018learning} have been applied to SSD
traces~\cite{chakraborttii2020learning} and to flash (a non-archival
workshop abstract~\cite{chakraborttii2022deep}). Online and streaming learners
adapt file-system prefetching as the workload
changes~\cite{savva2026adaptive,lucasfilho2023streaming}. Baleen learns
admission and prefetch for flash caches~\cite{wong2024baleen}, CAPSULE clusters
spatio-temporal locality for cloud block storage~\cite{ramadhan2026capsule},
KML runs models inside the kernel~\cite{akgun2021kml}, and CrossPrefetch
coordinates prefetching across the OS and the library~\cite{garg2024crossprefetch}.
Memory-side work learns remote-memory prefetch and page
placement~\cite{maruf2020leap,doudali2019kleio}. Such work usually compares a
learned model with one heuristic on a few traces; we fix the simple baselines
in advance and measure the other levers on the same traces.

\textbf{Reference points from cache theory.} Belady's MIN is the optimal
demand-paging replacement~\cite{belady1966study}; it does not bound a policy
that prefetches, and we never use it as such a bound. Stack-distance
analysis~\cite{mattson1970evaluation} is the classical evaluation tool, and
adaptive, FIFO-based and learned
eviction~\cite{megiddo2003arc,yang2023fifo,zhang2024sieve,yang2023glcache,
shi2019glider} change the replacement policy; we keep LRU fixed so that only
prefetching varies.

\textbf{Web prefetching and next-item models.} Predictive prefetching of web
pages~\cite{padmanabhan1996using}, its latency bounds~\cite{kroeger1997exploring}
and its cost over slow links~\cite{fan1999web} raised the same what and
how-much questions in the 1990s. Session-based next-item models range from
GRUs to self-attention~\cite{hidasi2016session,kang2018sasrec}; we test a GRU
as a ranker.

We do not propose a predictor; we measure, under pre-registered rules,
which lever moves the outcome, and check the simulators against a real disk
and real HTTP.
